\chapter{Protocolo de teste e regra de multiplicidade}
\label{cap:regra}

Cada análise declara a janela, o tamanho da amostra, a unidade que se repete
dentro dela e o teste aplicado. As declarações completas estão no
apêndice~\ref{ap:inventario}, e cada seção repete a sua na abertura.

\section{Janela, amostra e unidade de repetição}

As janelas de corte saem de uma regra única: percentil 80 do início e percentil
20 do fim do período de apuração de cada fonte. Elas diferem entre fontes porque
a apuração eólica começa em outubro de 2021 e a fotovoltaica em abril de 2024. A
análise da seção~\ref{sec:fig1} vem de extração anterior às demais, e por isso
sua taxa agregada difere em 0,5 ponto da extração final.

A unidade de repetição varia com a pergunta. Nas análises de despacho é a hora
do sistema, com $n$ da ordem de 21 mil. Nas análises de rede é o conjunto de
geração, com $n$ entre 59 e 120. A diferença de duas ordens de grandeza importa
para a leitura: um coeficiente pequeno é detectável no primeiro caso e não é no
segundo.

\section{Correção de multiplicidade}

A correção de Holm é aplicada dentro de cada análise, sobre as oito covariáveis
pré-declaradas dos capítulos de rede, e é declarada ausente onde não foi
aplicada. No estudo de caso pareado as cinco covariáveis testadas resultaram
todas nulas, e a correção só tornaria a conclusão mais conservadora.

Não há correção entre análises. Somados os testes que vão a arquivo, o relatório
reporta 49, dos quais 16 sobrevivem a Holm e 33 não. As oito covariáveis de rede
foram pré-declaradas e testadas na mesma ordem nas duas fontes de geração, e o
capítulo~\ref{cap:rede} submete as sobreviventes a três critérios adicionais. Um
limiar de multiplicidade em nível de relatório não é aplicado.

\begin{table}[t]
\centering
\tabfont
\begin{tabular}{@{}
    >{\rag\arraybackslash}p{2.9cm}
    >{\rag\arraybackslash}p{2.6cm}
    r@{}}
\toprule
\cab{Capítulo} & \cab{Onde entra} & \cab{Testes} \\
\midrule
Estudo de caso        & sem correção      &  5 \\
Detecção por satélite & duas comparações  &  2 \\
Despacho              & quintis de carga  &  5 \\
Decomposição          & sem correção      & 14 \\
Rede, solar           & Holm sobre oito   &  8 \\
Rede, eólica          & Holm sobre oito   &  8 \\
Robustez              & bootstrap, Fisher &  7 \\
\midrule
\cab{Total}           &                   & \cab{49} \\
\bottomrule
\end{tabular}
\caption{Distribuição dos testes que vão a arquivo.}
\label{tab:multiplicidade}
\notas{Dos 49 testes, 16 sobrevivem à correção de Holm aplicada dentro da
análise que os produz. Os 33 restantes são reportados sem correção, e a análise
que os contém declara essa ausência.}
\end{table}

\begin{objecao}
A correção de multiplicidade governa a taxa de erro dentro de cada análise, e
não a validade do resultado. Uma covariável que sobrevive a Holm pode ainda
assim não sobreviver a reamostragem. É o que a seção~\ref{sec:fig12} mede para
os quatro resultados de rede.
\end{objecao}
